\documentclass[10pt]{article}
\usepackage[margin=20mm]{geometry}
\usepackage{amsmath,amssymb,mathtools,fontspec,unicode-math,xurl,hyperref}
\IfFontExistsTF{DejaVu Serif}{\setmainfont{DejaVu Serif}}{\setmainfont{Latin Modern Roman}}
\setmathfont{latinmodern-math.otf}
\setlength{\parindent}{0pt}\setlength{\parskip}{6pt}\sloppy
\title{Scientific initial review of
Fractal uncertainty in higher dimensions}
\author{AI-assisted scientific review; human decision pending}
\date{}
\begin{document}
\maketitle
Abstract: This review concerns Alex Cohen’s arXiv:2305.05022v2 author manuscript, dated 5 October 2024, with 42 PDF pages. One supported Minor issue concerns the pointwise interpretation of a boundary-derivative lemma under a merely \(C^1\) hypothesis. An explicit compactly supported \(C^1\) example has infinite normal derivative at one boundary point. The subsequent core applications assume \(C^2\) or \(C^3\) regularity, so this example does not refute the fractal uncertainty theorem or the higher-dimensional multiplier construction. No Major issue was established in the inspected core. Several scale and normalization details are made explicit below, with local supporting arguments. All original pages and all supplied mathematical TeX/bibliography files were read; nine groups of reviewer-owned checks were executed. This is an initial review, not independent acceptance or certification of all external dependencies.\par
\section*{Identity, version, and scope}
The assigned paper is M25-0019, Alex Cohen, Fractal uncertainty in higher dimensions, DOI https://doi.org/10.4007/annals.2025.202.1.4. The actual scientific input is the manifest-bound arXiv:2305.05022v2 author manuscript. Its first page says 5 October 2024; the TeX filename contains 10-4-24 and the PDF creation metadata is later. Those are different metadata fields, not evidence of different mathematical versions.\par
All page references below mean the one-based pages of the 42-page acquired PDF; its printed page numbers agree. The journal Version of Record was not compared. Findings must not be attributed to that final version without a corresponding comparison.\par
The review covered the main statements, definitions and quantifiers, Sections 1–6, Appendix A, all 42 original page images, and the full 1782-line mathematical TeX file, 312-line bibliography output, and 446-line bibliography source. The three supplied illustrative figure assets were inspected as rendered in the main PDF on pages 2 and 33. No separate experimental data or numerical author program is present in this source bundle. No author code was executed. The review is analytic, with small symbolic and numerical cross-checks; it does not reproduce an empirical experiment.\par
\section*{Claims, assumptions, and checked dependence}
The Fourier convention is \[\widehat f(\xi)=\int_{\mathbb R^d}f(x)e^{-2\pi i x\cdot\xi}\,dx.\] Theorem 1.1, page 1, assumes \(X\subset[-1,1]^d\) is \(\nu\)-porous on balls from scales \(h\) to 1, and \(Y\subset[-h^{-1},h^{-1}]^d\) is \(\nu\)-porous on lines from scales 1 to \(h^{-1}\). Here \(0<h<1/100\), and the standing convention is \(0<\nu\leq1/3\). Porosity means that every tested ball or line segment has a point admitting a ball of radius \(\nu\) times its diameter or length that misses the set. The quantifier over all line segments is essential; ordinary ball porosity of both sets is insufficient.\par
The asserted conclusion is \[\operatorname{supp}\widehat f\subset Y
 \quad\Longrightarrow\quad
 \|f1_X\|_2\leq C h^\beta\|f\|_2,
 \qquad C,\beta>0\ \hbox{depending only on }\nu,d.\] The intermediate multiplier theorem, Theorem 1.4 on page 5, assumes a nonpositive weight \(\omega\), vanishing on \(B_2\), with \[|D^a\omega(x)|\leq C_{\rm reg}\langle x\rangle^{1-a}\quad(0\leq a\leq3),
 \qquad
 \int_0^\infty\frac{G^*(r)}{1+r^2}\,dr\leq C_{\rm gr},\] where \(G(x)=\int_{1/2}^2|\omega(sx)|\,ds\) and \(G^*(r)=\sup_{|x|=r}G(x)\). The positive constants in upper bounds may always be enlarged, including in the zero-weight case. The theorem supplies a nonzero \(L^2\) multiplier with Fourier support in \(B_\sigma\), a uniform lower bound near zero, and exponential decay controlled by \(\sigma\omega\).\par
The logical dependence inspected was: Proposition 3.1 and Lemmas 4.1–4.5 give the plurisubharmonic extension in Proposition 2.2; the weighted \(\bar\partial\) argument in Section 5 gives Proposition 2.3; their combination gives Theorem 1.4; line porosity gives the damping functions in Proposition 1.7; the stated Han–Schlag conditional theorem, with the conversions in Appendix A.1, yields Theorem 1.1. The imported weighted \(\bar\partial\) and Paley–Wiener theorems were treated as established external results with their stated hypotheses. Their full original proofs were not independently reconstructed.\par
\subsection*{Extension and real-boundary positivity}
For the operator \[E\omega(x+iy)=\frac1\pi\int_\mathbb R\frac{\omega(x+ty)}{1+t^2}\,dt,\] the compact-support calculation on pages 17–19 is consistent: \[\partial\bar\partial E\omega(x+iy)
 =\frac1{4\pi}\int_\mathbb RD^2\omega(x+ty)\,dt.\] The mixed imaginary coefficients cancel by equality of mixed derivatives. Changing variables gives the factor \(1/|y|\). The X-ray transform is constant along its integration line, so its Hessian annihilates the longitudinal direction. For \(v=v_\perp+r\widehat y\) this gives the transverse formula (3.15). Meanwhile, \[\partial\bar\partial|y|=\frac{I-\widehat y\widehat y^{\,T}}{4|y|}.\] Consequently the line-integral lower bound in Proposition 3.1 offsets exactly the possible negative transverse Hessian. No missing \(4\), \(\pi\), or \(|y|\) factor was found here.\par
At a real point, pages 20–21 use \(C>C_1\) first. The positive term \(\varepsilon\,\operatorname{avg}_\theta|\operatorname{Im}(e^{i\theta}v)|\) is at least \(\varepsilon|v|/2\) and dominates the negative quadratic Taylor remainder for small \(|v|\). Passing \(C\) down to its limiting value is legitimate by local uniform convergence. This step uses \(C^2\) regularity, not merely the hypothesis challenged in Problem 1.\par
For noncompact weights, the uniform tail estimate on page 16 includes both a term of size \(R^{-1/2}C_{\rm gr}\) and the integrable tail \(M\int_{\sqrt R}^{\infty}G^*(t)/(1+t^2)\,dt\). It therefore gives convergence uniformly on compact subsets, including as \(y\) tends to zero. The same domination supports Lemma 3.5 for the specific nonpositive partial sums constructed later.\par
\subsection*{Dyadic modification and cancellation}
The integration-by-parts sign in Lemma 4.1, page 22, was checked: \[\int_0^\infty D^2_{vv}f(t\widehat y)\,dt
 =\pi_{\mathbb S^{d-1}}f(\widehat y)
  +\left.\frac{d^2}{d\theta^2}
    \pi_{\mathbb S^{d-1}}f(\widehat y\cos\theta+v\sin\theta)
   \right|_{\theta=0}.\] Compact support away from the origin removes the boundary term. Making each large dyadic piece have constant spherical projection \(q_k\) leaves the line integral equal to \(2q_k\) on the full line. The correction \(p_k^{-1}\psi_k(q_k-\pi\omega_k)\) is nonpositive, so the new weight stays below the old one. The summability of \(|q_k|\), together with the third derivative bound, controls the large-scale translation errors in Lemmas 4.4–4.5. The low, middle, and high scales have respectively summable geometric bounds, a bounded number of terms, and the summable \(|q_k|\) terms.\par
The audit retains the nonpositivity of the constructed pieces in (4.13). Reading Lemma 4.3 as a theorem about arbitrary signed pieces solely from the abbreviated list (ii)–(v) would discard information used in the absolute-integral bounds. The actual constructed pieces have that information; no defect in their use was established. Likewise the displayed projection normalization \(p_k\asymp2^{-k}\) is used for \(k\geq1\), not for the partition piece meeting the origin.\par
\subsection*{The weighted \(\bar\partial\) step}
Pages 28–31 use \[\varphi=2u+20d\log|z|_\infty+\rho u_0
              +\frac{\rho}{2}(\langle y\rangle-1).\] The singular logarithm is plurisubharmonic; the source expressly permits distributional Hessians in the invoked theorem. The positive correction has \[\partial\bar\partial\langle y\rangle
 =\frac{(1+|y|^2)I-yy^T}{4(1+|y|^2)^{3/2}},
 \qquad \lambda_{\min}=\frac1{4\langle y\rangle^3}.\] Thus the stated lower bound \(\kappa=\rho/(8\langle y\rangle^3)\) is correct. The cutoff derivative is supported away from zero, where the weight is bounded below. The solution \(g\) is holomorphic near zero, and the weighted submean estimate gives \[|g(x)|\leq C_d\rho^{-1/2}|x|^{9d}e^{3\rho|x|}.\] For \(|x|\leq c_d\min(\rho,\rho^{-1})\), the right side is uniformly small after choosing \(c_d\) sufficiently small. For \(\rho\leq1\) the power is \(\rho^{9d-1/2}\); for \(\rho\geq1\) it is \(\rho^{-9d-1/2}\), while the exponential is bounded in both regimes.\par
The additional negative weight at infinity absorbs logarithmic growth and makes the real restriction square integrable. The complex bound \(Ae^{2\rho|y|}\) gives Fourier support in \(B_{\rho/\pi}\) under this Fourier convention. The rescaling on page 14 then gives the announced support in \(B_\sigma\). Appendix A.3 is explicitly a sketch; its contour-shift sketch is not a substitute here for a full Paley–Wiener proof.\par
\subsection*{Line porosity and the exponent}
In Section 6, \(W_k=2^k/k^s\) and the weight height is \(2^k/k^\alpha\). Its derivatives through order 3 obey the required symbol bounds when \(3s<\alpha\). Restricting thickened porous sets to a line and using Appendix A.2 yields a gain \(k^{-s\gamma}\). Accordingly, \[G^*(r)\lesssim \frac{r}{(\log(2+r))^{\alpha+s\gamma}},\] and integrability requires \(\alpha+s\gamma>1\). After reducing \(\gamma\) if needed to \(0<\gamma\leq1\), the choice \(s=1/5\), \(\alpha=1-\gamma/10\) gives \[\alpha-3s=2/5-\gamma/10\geq3/10,\qquad
 \alpha+s\gamma-1=\gamma/10>0.\] These are compatible strict margins. The porosity counting identity \((L^d-1)^nL^{-nd}=(1-L^{-d})^n\) and the exponent \(-\log(1-L^{-d})/\log L\) also agree with the volume argument.\par
\section*{Problem 1: pointwise normal derivative under a \(C^1\) hypothesis}
\subsection*{Source claim and exact location}
Judgment: Minor, for the classical pointwise reading. Lemma 3.3, PDF page 17, equation (3.9), asserts for \(\omega\in C^1_0(\mathbb R)\) and its bounded harmonic extension to the upper half-plane that \[\partial_yu(x+i0)=H[-\omega'].\] The TeX label is `lem:normal\_derivative\_harmonic\_extension`. The proof ends by passing through the Hilbert transform from \(C^1\) convergence of the boundary approximations. Such convergence does not imply finite pointwise convergence of the Hilbert transforms of their continuous derivatives.\par
\subsection*{Explicit admissible example and complete check}
Put \(r_0=1/8\) and \(r_1=1/4\). Define an even cutoff \(\chi\) by \[\chi(x)=
 \begin{cases}
 1,& |x|\leq r_0,\\
 1-3s^2+2s^3,\quad s=(|x|-r_0)/(r_1-r_0),
       &r_0<|x|<r_1,\\
 0,&|x|\geq r_1.
 \end{cases}\] The polynomial has zero derivative at both endpoints and takes values in \([0,1]\). Thus \(\chi\) is compactly supported and \(C^1\). Set \[\omega(0)=0,\qquad
 \omega(x)=\frac{\chi(x)|x|}{\log(e/|x|)}
   \quad(0<|x|<r_1),\qquad
 \omega(x)=0\quad(|x|\geq r_1).\] Near zero, with \(L(t)=\log(e/t)\), \[\frac{\omega(x)}x\longrightarrow0,\qquad
 \omega'(x)=\operatorname{sgn}(x)
       \bigl(L(|x|)^{-1}+L(|x|)^{-2}\bigr)\longrightarrow0.\] The cutoff endpoints have matching values and derivatives. Hence \(\omega\in C^1_0(\mathbb R)\) exactly as required, is real-valued, bounded and nonnegative, and \(\omega'(0)=0\). Lemma 3.3 does not require a sign; replacing \(\omega\) by \(-\omega\) also gives a nonpositive example if a weight convention is imposed.\par
Equation reference: eq:quotient. Its Poisson integral is the bounded harmonic extension required by the lemma. Positivity and the unit mass of the Poisson kernel give \(0\leq u\leq\|\omega\|_\infty\). For \(0<y<r_0\), \[\begin{aligned}
 \frac{u(iy)-u(0)}{y}
 &=\frac2\pi\int_0^{r_1}
          \frac{\omega(t)}{t^2+y^2}\,dt\\
 &\geq\frac1\pi\int_y^{r_0}\frac{dt}{tL(t)}
 =\frac1\pi\log\frac{L(y)}{L(r_0)}
 \longrightarrow+\infty. \end{aligned}\] This proves that the finite classical boundary normal derivative does not exist.\par
Equation reference: eq:trace. Even if the notation in (3.9) is read as a limit of the interior derivatives, the same failure occurs. Differentiation for \(y>0\) gives \[\partial_yu(iy)=\frac2\pi\int_0^{r_1}
       \omega(t)\frac{t^2-y^2}{(t^2+y^2)^2}\,dt.\] The absolute value of the negative contribution from \(0<t<y\) is at most \(1/(\pi L(y))\). For \(2y\leq t\leq r_0\), the kernel is at least \(12/(25t^2)\); all other terms with \(t\geq y\) are nonnegative. Therefore, for \(2y<r_0\), \[\partial_yu(iy)\geq
 \frac{24}{25\pi}\log\frac{L(2y)}{L(r_0)}
       -\frac1{\pi L(y)}
 \longrightarrow+\infty. \] There is also no finite principal value on the stated right side at zero. The inner contribution is \[H[-\omega'](0)
 =\frac2\pi\int_0^{r_0}
    \frac{L(t)^{-1}+L(t)^{-2}}{t}\,dt
       +\hbox{a finite outer contribution}
 =+\infty.\] The divergence of \(\int_0^{r_0}dt/(t\log(e/t))\) is exact and is not a numerical inference.\par
\subsection*{Counterevidence, impact, and verified correction}
The distributional identity is valid; so is an appropriate \(L^2\) boundary-trace formulation. Since \(\omega'\) is compactly supported and in \(L^2\), Hilbert-transform continuity in \(L^2\) supplies that weaker interpretation. The lemma need not be read as claiming a new failure of that standard operator identity. Its pointwise notation and the displayed limiting argument should specify the intended topology.\par
A sufficient classical correction is to assume \(\omega\in C_c^{1,\alpha}(\mathbb R)\) for some \(0<\alpha\leq1\), or more generally a Dini modulus for \(\omega'\). In the symmetric principal value, the local numerator is bounded by a constant times \(t^\alpha\), so the integrand is integrable as \(t^{\alpha-1}\). The tails are finite. The resulting continuous boundary function is recovered by Poisson approximation, giving the finite classical normal derivative and the identity. In particular, \(C_c^2\) is sufficient.\par
This defect is peripheral to the inspected core. Lemma 3.4 on the same page assumes \(C^2_0(\mathbb R^d)\); its one-dimensional application on pages 17–18 uses \(C^2_0(\mathbb R)\) explicitly. Proposition 3.1 also assumes \(C^2\), the modified weights have three controlled derivatives, and Lemma A.9 assumes \(C^2_0\). Those uses meet the sufficient corrected hypothesis. The example verifies an overstated pointwise regularity range, not a counterexample to Theorems 1.1 or 1.4. A wording clarification in a distributional formulation, or the stated strengthening for a classical formulation, resolves the identified issue. Author adoption and final-version status remain unknown.\par
\section*{Scope clarifications with local verification}
The following are recorded to make the checked argument reproducible. They are not additional supported Major findings and are not counted as separate typographical complaints.\par
\subsection*{The final Schwartz factor, pages 34–35}
The text introduces a scaled factor and then suppresses its scale in the product notation. Its normalization should ensure a nonzero value at zero, not merely an integral normalization for the untransformed factor. A fully specified choice is to take \(\eta\in C_c^\infty(B_1)\), \(\eta\geq0\), \(\int\eta=1\), put \(\Phi=\check\eta\), and use \(\Phi_\delta(x)=\Phi(\delta x)\) with \(\delta=c_d\sigma\leq\sigma/2\). Then \(\Phi(0)=1\), \(|\Phi|\leq1\), and \(|\Phi_\delta|\geq1/2\) on a sufficiently small ball of radius \(c'_d\sigma\), since \(0<\sigma<1\).\par
If \(\operatorname{supp}\widehat f\subset B_{\sigma/2}\), the product has Fourier support in \(B_\sigma\). The near-origin lower bound is at least \(1/4\), which is enough after changing constants. For each fixed decay order, \(|\Phi_\delta(x)|\leq C_{d,M}\delta^{-M}\langle x\rangle^{-M}\). Dividing the product by a sufficiently large constant times a negative power of \(\sigma\) gives the prescribed global \(\langle x\rangle^{-d}\) upper bound, retains exponential decay on \(Y\), and gives a positive lower \(L^2(B_1)\) bound that is a positive power of \(\sigma\). This verifies the intended damping construction without asserting the unscaled literal product has the smaller Fourier support.\par
\subsection*{The upper end of the line scales, pages 33–34}
The segment length \(2^k/\sqrt d\) may exceed \(N=h^{-1}\) for the outermost nonempty annuli. A direct way to stay within the stated porosity range is to use \[R_k=\min(2^k/\sqrt d,N/2).\] Since a nonempty annulus has \(2^k\leq3\sqrt d\,N\), this length is comparable to \(2^k\) with constants depending only on \(d\). Cover the intersection with the line by a bounded number, depending on \(d\), of such segments. If the lower porous scale \(a_k=4W_k\sqrt d/\nu\) is below \(R_k\), Corollary A.8 gives exactly the same power \(k^{-s\gamma}\). If \(a_k\geq R_k\), then \(k^s\leq24d/\nu\); the trivial length bound is therefore at most a constant depending on \(\nu,d\) times \(2^k k^{-s\gamma}\). Thus (6.10)–(6.12) retain their asserted parameter dependence. This is a verified local scale adjustment, not a demonstrated loss of the main estimate.\par
\subsection*{Translations and bounded small scales, page 35}
Writing \(R=sN\), the translated set is contained in \([-2R-9,2R+9]^d\). The box \([-3R,3R]^d\) required for a direct use of Proposition 1.7 contains this when \(R\geq9\). For \(R\geq102\) the small-\(h\) convention is also harmless. For \(1<R<102\), all these sets lie in the fixed box \([-213,213]^d\). A fixed bandlimited Schwartz factor with nonzero value at zero, scaled to the prescribed positive \(c_1=\nu/(20\sqrt d)\) and multiplied by a sufficiently small positive constant, meets both the polynomial bound globally and the exponential bound on this fixed box. Its local \(L^2\) mass remains a positive constant depending only on \(\nu,d\). Hence the bounded-scale cases do not require porosity at an unavailable larger scale.\par
\subsection*{The endpoint \(s=1\) and geometric slack, Appendix A.1}
The imported iteration includes depth \(n=0\), whereas Theorem 1.6 formulates its damping assumption for \(s<1\). The endpoint follows from compactness: choose \(s_j\uparrow1\) with uniform damping parameters. The common bound \(\langle x\rangle^{-d}\) gives a common \(L^2\) bound. Common compact Fourier support gives local equicontinuity by Cauchy–Schwarz on the Fourier side. A subsequence converges locally uniformly and weakly in \(L^2\); its Fourier support is unchanged, its local lower norm survives, and the pointwise decay passes along points of the moving sets \(s_jY+[-4,4]^d+\eta\). The two-unit slack between the translation ranges in (A.5) accommodates each fixed endpoint translation for large \(j\).\par
For strict ball-scale endpoints and cube-boundary containment one can decrease the working porosity constant and use a larger integer \(L\), for example \(\lceil4\sqrt d/\nu\rceil\). At every required depth \(L^{n+1}\leq\lceil h^{-1}\rceil\), a ball of diameter half the parent side lies in the allowed scale range; the smaller child cube then lies strictly inside a porosity hole. The chosen Fourier radius still fits since \(\nu/(20\sqrt d)<1/(2L)\) for \(0<\nu\leq1/3\). Constants in the final theorem are unspecified functions of \(\nu,d\), so this conservative slack preserves its assertion.\par
\section*{Overall impact, references, and remaining uncertainty}
Item | Judgment | Repair and remaining work\par
Lemma 3.3 | Minor for finite pointwise interpretation | \(C^{1,\alpha}_c\) correction proved; weak identity retained. Adoption unknown.\par
Core multiplier and FUP chain | No Issue Found within inspected scope | Main calculations checked; no independent certification of every imported theorem.\par
Scale/normalization details | Scope clarification | Local constructions and parameter slack supplied; not additional Major findings.\par
Final journal version | Not established | Compare corresponding final text before transferring any finding.\par
The external check of Han–Schlag concerned their arXiv:1805.04994v2, Definition 4.1 and Theorem 5.1, on printed pages 20 and 24. Their support, translation, and discrete-scale conditions match the dependency used here after the stated norm conversions and endpoint clarification. This selected check is not a full rereading of that paper or a comparison with its published version. The \(\ell^1/\ell^2\) comparison in the present Appendix A.1 is consistent with reducing the decay constant by \(\sqrt d\).\par
The arXiv metadata for the assigned paper identifies v2 as its displayed latest revision at the time of this check. Bounded title/erratum searching did not establish a published correction. This is not evidence that none exists. Dyatlov’s arXiv:2305.17857 was located at metadata level only; its proof was not used as a substitute for the present review. No fresh empirical experiment is required by the identified analytic issue.\par
9 A. Cohen, Fractal uncertainty in higher dimensions, arXiv:2305.05022v2, 5 October 2024. Manifest-bound 42-page author manuscript, fully read. https://arxiv.org/abs/2305.05022v2. Journal identifier: https://doi.org/10.4007/annals.2025.202.1.4; final journal text not compared.\par
R. Han and W. Schlag, A higher dimensional Bourgain–Dyatlov fractal uncertainty principle, arXiv:1805.04994v2, 22 May 2018. Selected dependency statements consulted, not full external proof reproduction. https://arxiv.org/abs/1805.04994v2.\par
S. Dyatlov, Fractal uncertainty in higher dimensions: notes on Cohen’s paper, arXiv:2305.17857v1, 29 May 2023. Metadata only; not a consulted proof. https://arxiv.org/abs/2305.17857.\par
\section*{Supporting checks}
Process 23912 exited with code 0. The environment reported Python 3.13.5, SymPy 1.13.3, and mpmath 1.3.0. The results are in `reviewer\_checks\_results.json`. Nine groups passed: the local derivative of the counterexample; five Poisson-integral illustrations; eight extension-integrand identities; a compact \(C^3\) test of Lemma 4.1; the strict-PSH Hessian in dimension 3; 100 rational exponent cases; 192 exact porosity-recursion cases; 68 illustrative nonvanishing-radius cases; and the elementary Fourier-radius slack.\par
The symbolic compact test of Lemma 4.1 used \(\psi(r)=(r-1)^4(2-r)^4\) on \([1,2]\), extended by zero, with angular factor \(\cos^2\theta\). It gave \[q=\frac{2329}{35}-96\log2,\qquad
 \int_1^2\left(\frac{\psi'(r)}r-\frac{2\psi(r)}{r^2}\right)\,dr=-q,\] in agreement with the spherical second derivative. These finite checks supplement, and do not replace, the general arguments above.\par
For the counterexample, numerical integration over \(y\leq t\leq r_0\) gives approximately \(0.337384\), \(0.740060\), \(1.160663\), \(1.591134\), and \(2.026884\) for \(y=10^{-2},10^{-4},10^{-8},10^{-16},10^{-32}\). These are truncated positive contributions to the normal quotient, not full Poisson solutions. Divergence follows from the analytic lower bound [eq:quotient] (\#eq:quotient); no finite table proves it.\par
Static source auditing read and hashed the supplied TeX and bibliography files and checked local labels/references. It did not execute TeX, author macros, author scripts, or author experiments. The main PDF was extracted with Poppler and all 42 rendered images were inspected. Rendering emitted a missing-display-font warning for Symbol; the reviewed pages, including the decisive formulas and all three figures, were readable. No missing mathematical glyph affecting this finding was observed.\par
No `pdflatex`, `xelatex`, `lualatex`, or `tectonic` was found on the checked PATH or recorded standard locations. No installation was attempted. Static environment-balance checking is not a substitute. The acquired original PDF and its page images are evidence, not a compiled report.\par
\section*{Supporting checks}
Task token:\par
The PDF SHA-256 is\par
> 6424ce6c3926bce57c98c99492d38bcf5bfc4cf30b2c59592f15e43c779aba00\par
The mathematical TeX SHA-256 is\par
> 4e5198451e311c6932427b4d65d892b59d37342b39db672821c2c5fd4dd52164\par
The manifest SHA-256 is\par
> 7b16c2183a7be4810adbac8353079f12936acd865dfe91ec32830255d8bff245\par
All 12 manifest-listed files were verified against their supplied hashes.\par
This assistant did not personally inspect that dispatcher-held browser tab. No model switch was performed, and no independent evidence of a non-6Pro backend route was observed.\par
All scientific analysis, source inspection, mathematics, reviewer-owned computation and original-page inspection described here were performed by this assistant. No scientific subagent was used. No human approval, author adoption, or journal endorsement is represented. Files were saved only under the assigned Windows staging task. No final archive, cloud upload, author contact, installation, payment, or persistent heartbeat was created.\par
\section*{Compilation continuation (26 October 2026)}
Scientific conclusions, pending questions, source-version limits and the absence of independent acceptance remain unchanged.\par
\end{document}
